\subsection{トピックと参考文献の対応}
この表は、SWEBOK が示す「トピック対参考資料の行列」を、学習用に簡略化したものである。詳しく学ぶときは、各トピックに対応する文献を読む。

\par\smallskip\noindent\refstepcounter{table}\label{tab:ch09-08}表 \thetable: トピックと参考文献の対応におけるトピック・主な参考資料\par\smallskip
表 \thetable は、トピックと参考文献の対応におけるトピック・主な参考資料を比較・確認しやすい形で整理したものである。\par\smallskip
\begin{longtable}{>{\raggedright\arraybackslash}p{0.25\linewidth}>{\raggedright\arraybackslash}p{0.70\linewidth}}
\toprule
トピック & 主な参考資料 \\
\midrule
\endfirsthead
\toprule
トピック & 主な参考資料 \\
\midrule
\endhead
1 開始と範囲定義 & Fairley, Managing and Leading Software Projects; Sommerville, Software Engineering \\
1.1 要求の決定と交渉 & Fairley; Software Requirements KA \\
1.2 実現可能性分析 & Sommerville \\
1.3 要求のレビューと改訂 & Fairley; Software Configuration Management KA \\
2 プロジェクト計画 & Fairley; Boehm and Turner; Sommerville \\
2.1 プロセス計画 & Fairley; Boehm and Turner \\
2.2 成果物の決定 & Fairley \\
2.3 工数・日程・費用見積り & Fairley; Software Engineering Economics KA \\
2.4 資源割当て & Fairley \\
2.5 リスク管理 & Fairley; Boehm and Turner \\
2.6 品質管理 & Fairley; Sommerville; Software Quality KA \\
2.7 計画管理 & Fairley \\
3 プロジェクト実行 & Fairley; Sommerville \\
3.2 取得と供給者契約管理 & Fairley \\
3.4 監視プロセス & Fairley \\
3.5 制御プロセス & Fairley \\
4 レビューと評価 & Fairley; Sommerville \\
5 終結 & PMBOK Guide; Software Extension to PMBOK Guide \\
6 測定 & Practical Software Measurement; ISO/IEC/IEEE 15939 \\
7 管理ツール & Fairley; Practical Software Measurement \\
\bottomrule
\end{longtable}
